\subsection{Supervision, Scoring and Post-hoc Deferral}
\label{sec:ablation}
\noindent\begin{minipage}{\columnwidth}
\centering
\captionof{table}{Full-budget means (\%). Avg.: equal-task mean (five-seed SD). $\Delta$: gain over $\mathrm{LR}_{PH}$ (pp). Shared weights for $\mathrm{LR}_{4S}$. Earlier selection settings in supplement. CS: ConstructionSite.}
\label{tab:objective}
\fontsize{8}{9.5}\selectfont
\setlength{\tabcolsep}{1.5pt}
\begin{tabular}{@{}llrrrrrr@{}}
\toprule
Variant & Score & CUB & gRef & NLVR & CS & Avg. & $\Delta$ \\
\midrule
$\mathrm{LR}_{PH}$ & $r$ & 65.86 & 59.37 & 87.33 & 21.24 & 58.45 (0.07) & -- \\
\midrule
3-state & $p_R-p_H$ & 71.27 & 60.13 & 86.34 & 21.77 & 59.88 (0.06) & \textcolor{upgreen}{+1.43} \\
Gain & $\widehat\Delta$ & 71.01 & 59.97 & 86.29 & 21.61 & 59.72 (0.04) & \textcolor{upgreen}{+1.27} \\
\midrule
$\mathrm{LR}_{4S}$ & $p_R-p_H$ & 71.30 & 60.10 & 86.43 & 21.69 & 59.88 (0.06) & \textcolor{upgreen}{+1.43} \\
$\mathrm{LR}_{4S}$ & $\log(p_R/p_H)$ & 67.80 & 59.84 & 87.36 & 21.11 & 59.03 (0.09) & \textcolor{upgreen}{+0.58} \\
$\mathrm{LR}_{4S}$ & $r_{4S}$ & 67.63 & 59.90 & 87.36 & 21.30 & 59.05 (0.10) & \textcolor{upgreen}{+0.60} \\
\rowcolor{learnfill}
$\mathrm{LR}_{4S}$ & $s_\alpha$ & 71.37 & 60.66 & 87.36 & 21.63 & 60.25 (0.06) & \textcolor{upgreen}{+1.81} \\
\bottomrule
\end{tabular}

\end{minipage}\par\vspace{4pt}
Table~\ref{tab:objective} compares training objectives and scores. Merging the two unchanged outcomes into one target or directly predicting the gain gives curves close to the four-state probability difference. The supplement gives their losses and training settings.

Changing only the score has a larger effect on some tasks. Here $r_{4S}=\log[(p_R+p_B)/(p_H+p_B)]$ estimates the optimal score of the post-hoc loss~\cite{posthocdefer}, whereas $\mathrm{LR}_{PH}$ learns its score directly. On NLVR2, this change largely closes the middle-to-high-budget gap to post-hoc deferral, but it hurts CUB. The normalized score has the highest task average, while no score is best on every task.
